\UseRawInputEncoding
\documentclass[12pt,openright,twoside,a4paper,english,brazilian,hidelinks]{abntex2}

% Pacotes essenciais
\usepackage[utf8]{inputenc}        % UTF-8
\usepackage[T1]{fontenc}           % Fontes T1
\usepackage{lmodern}               % Fontes Latin Modern

\usepackage{xcolor}                % Cores
\usepackage{graphicx}              % Imagens
\usepackage{geometry}              % Ajustes pontuais de margem
\usepackage{amsmath}               % Matemática
\usepackage{listings}              % Código fonte
\usepackage{xspace}                % Espaçamento automático
\usepackage{float}                 % Posicionamento de figuras
\usepackage{pdflscape}             % Páginas em paisagem no PDF
\usepackage{hyperref}   % Hiperlinks sem caixas

% Diretório de gráficos e diagramas
\graphicspath{{design_template/diagramas/}}

\newcommand{\diagramPortraitPage}[3]{%
    \clearpage
    \thispagestyle{empty}%
    \vspace*{\fill}
    \begin{center}
    \textbf{#1}\par\vspace{0.5cm}
    \includegraphics[width=\textwidth,height=0.80\textheight,keepaspectratio]{#2}
    \refstepcounter{figure}\par\vspace{0.35cm}
    \textbf{Figura \thefigure{} -- #3}
    \end{center}
    \vspace*{\fill}
    \clearpage
}

\newcommand{\diagramLandscapePage}[3]{%
    \clearpage
    \begin{landscape}
    \thispagestyle{empty}%
    \vspace*{\fill}
    \begin{center}
    \textbf{#1}\par\vspace{0.5cm}
    \includegraphics[width=\linewidth,height=0.80\textheight,keepaspectratio]{#2}
    \refstepcounter{figure}\par\vspace{0.35cm}
    \textbf{Figura \thefigure{} -- #3}
    \end{center}
    \vspace*{\fill}
    \end{landscape}
    \clearpage
}

\newcommand{\diagramWideLandscapePage}[3]{%
    \clearpage
    \newgeometry{left=1cm,right=1cm,top=1cm,bottom=1cm}
    \begin{landscape}
    \thispagestyle{empty}%
    \vspace*{\fill}
    \begin{center}
    \textbf{#1}\par\vspace{0.5cm}
    \includegraphics[width=0.98\linewidth,keepaspectratio]{#2}
    \refstepcounter{figure}\par\vspace{0.35cm}
    \textbf{Figura \thefigure{} -- #3}
    \end{center}
    \vspace*{\fill}
    \end{landscape}
    \restoregeometry
    \clearpage
}

% Configuração de código fonte
\lstset{
    basicstyle=\ttfamily\small,
    keywordstyle=\color{blue},
    commentstyle=\color{gray},
    stringstyle=\color{red},
    breaklines=true,
    showstringspaces=false,
    tabsize=4,
    literate={é}{{\'e}}1,
    language=Java,
    frame=single,
    rulecolor=\color{black!30}
}

% Espaçamento 1.5
\renewcommand{\baselinestretch}{1.5}

% ============================================================================
% INFORMAÇÕES DE IDENTIFICAÇÃO PESSOAL (PREENCHER CONFORME NECESSÁRIO)
% ============================================================================

\newcommand{\autorNome}{JAFTE CARNEIRO FAGUNDES DA SILVA}
\renewcommand{\instituicao}{ESCOLA POLITÉCNICA}
\newcommand{\universidade}{[PONTIFÍCIA UNIVERSIDADE CATÓLICA DO PARANÁ]}
\newcommand{\curso}{CIÊNCIA DA COMPUTAÇÃO}
\newcommand{\disciplina}{RESOLUÇÃO DE PROBLEMAS COM GRAFOS}
\newcommand{\professsor}{FABRÍCIO ENEMBRECK }
\newcommand{\semestre}{2026.1}
\newcommand{\dataprojeto}{maio de 2026}

% ============================================================================
% CONFIGURAÇÃO DO DOCUMENTO
% ============================================================================

\titulo{Analisador de Contatos Enron: Implementação de Algoritmos de Grafos\\em Java}

\autor{\autorNome}



\local{Curitiba}

\data{\dataprojeto}

\preambulo{
    Relatório de Aprendizado apresentado como trabalho prático da disciplina
    \textit{\disciplina{}} no curso de \curso{} da \universidade{}.

    \vspace*{\fill}

    \noindent\textbf{Orientador:} \professsor{}

    \vspace*{\fill}
}

\makeindex

% ============================================================================
% DOCUMENTO
% ============================================================================

\begin{document}

    \selectlanguage{brazilian}

    \frenchspacing

    \imprimircapa

% FOLHA DE ROSTO
    \begin{folhaderosto}
        \begin{center}
            \textbf{\autorNome} \\
        \end{center}

        \vspace*{3cm}

        \begin{center}
            \textbf{\imprimirtitulo}
        \end{center}

        \vspace*{3cm}

        \begin{flushleft}
            \preambulo{}
        \end{flushleft}

        \vspace*{\fill}

        \begin{center}
            Data: \imprimirdata
        \end{center}
    \end{folhaderosto}

% SUMÁRIO
    \pdfbookmark[0]{\contentsname}{toc}
    \tableofcontents*
    \cleardoublepage

% ============================================================================
% RESUMO
% ============================================================================

    \begin{resumo}
        Este relatório documenta o desenvolvimento de um analisador de contatos baseado no
        \textit{Enron Email Dataset} público, implementado em Java 17 como exercício prático
        da disciplina \disciplina{}. O projeto implementa, do zero, algoritmos clássicos de
        teoria dos grafos --- busca em profundidade (DFS), busca em largura (BFS), cálculo de
        distância e caminho crítico via Dijkstra adaptado --- aplicados a um grafo direcionado,
        ponderado e rotulado que modela a rede de comunicação entre funcionários da empresa.

        A arquitetura segue rigorosamente os princípios SOLID de engenharia de software,
        com separação clara de responsabilidades em 7 pacotes e 29 classes. Todos os algoritmos
        tratam ciclos explicitamente via conjuntos de visitados. A interface gráfica, construída
        com Java Swing, oferece 4 abas de análise interativa com visualização de grafos via
        GraphStream.

        \textbf{Palavras-chave:} Grafos. DFS. BFS. Dijkstra. Java. Enron. Estruturas de Dados.
    \end{resumo}

% ============================================================================
% INTRODUÇÃO
% ============================================================================

    \chapter{Introdução}

    A teoria dos grafos é uma das áreas fundamentais da ciência da computação, com aplicações
    práticas em redes sociais, sistemas de roteamento, análise de dependências e muitas outras
    áreas. Neste trabalho, desenvolvemos um analisador de contatos baseado no \textit{Enron
    Email Dataset}, um conjunto de dados públicos contendo mais de 600 mil mensagens de e-mail
    de aproximadamente 150 funcionários da empresa Enron, coletadas entre 1998 e 2002.

    O objetivo do projeto foi consolidar conhecimentos sobre estruturas de dados e algoritmos
    de grafos através de uma implementação prática completa, desde a representação do grafo até
    a execução de consultas complexas. Para isso, implementamos:

    \begin{enumerate}
        \item Um grafo direcionado, ponderado e rotulado em memória
        \item Busca em profundidade (DFS) iterativa
        \item Busca em largura (BFS) com garantia de caminho mínimo
        \item Cálculo de distância D entre nós
        \item Caminho crítico aproximado via Dijkstra adaptado com custo inverso
    \end{enumerate}

    Todos os algoritmos foram implementados \textbf{do zero}, sem uso de bibliotecas prontas
    de estruturas de dados ou algoritmos, utilizando apenas Java Collections (HashMap, HashSet,
    ArrayList, ArrayDeque) e a biblioteca padrão do Java 17.

    ---

    \section{Motivação}

    A análise de redes de comunicação é relevante para entender padrões de colaboração,
    identificar pessoas-chave em organizações e detectar caminhos críticos de informação.
    O \textit{Enron Email Dataset} representa um caso de estudo real, permitindo explorar
    estes tópicos com dados autênticos.

    ---

% ============================================================================
% REFERENCIAL TEÓRICO
% ============================================================================

    \chapter{Referencial Teórico}

    \section{Grafos Direcionados, Ponderados e Rotulados}

    Um grafo é uma estrutura matemática composta por um conjunto de vértices (ou nós) e um
    conjunto de arestas (ou arcos) que conectam pares de vértices. Neste projeto, utilizamos:

    \begin{itemize}
        \item \textbf{Direcionado}: cada aresta tem uma direção (origem → destino). A aresta
        $A \to B$ é distinta de $B \to A$.
        \item \textbf{Ponderado}: cada aresta tem um peso (neste caso, a frequência de
        mensagens enviadas de origem para destino).
        \item \textbf{Rotulado}: cada vértice tem um rótulo (o endereço de e-mail do indivíduo).
    \end{itemize}

    \textbf{Representação interna:} lista de adjacência via mapa de mapas:

    \begin{lstlisting}[language=Java]
Map<Vertex, Map<Vertex, Edge>> adjacency
// adjacency.get(origin).get(destination) = Edge com peso
    \end{lstlisting}

    Esta representação oferece $O(1)$ lookup de arestas e $O(k)$ enumeração de vizinhos,
    onde $k$ é o grau do vértice.

    ---

    \section{Busca em Profundidade (DFS)}

    A busca em profundidade é um algoritmo de traversal que explora o grafo ``o máximo
    possível'' em cada ramo antes de retroceder.

    \textbf{Características:}
    \begin{itemize}
        \item Usa uma \textbf{pilha} (LIFO)
        \item Complexidade: $O(V + E)$
        \item Não garante caminho mínimo
        \item Útil para detectar componentes conexas e ciclos
    \end{itemize}

    \textbf{Pseudocódigo:}

    \begin{lstlisting}[language=Java]
DFS(origem, destino):
    stack.push(origem)
    visited = {}
    predecessor = {}

    while stack not empty:
        current = stack.pop()
        if current in visited: continue
        visited.add(current)

        if current == destino:
            return reconstructPath(predecessor, origem, destino)

        for each neighbor in adjacency[current]:
            if neighbor not in visited:
                predecessor[neighbor] = current
                stack.push(neighbor)

    return empty  // nenhum caminho encontrado
    \end{lstlisting}

    ---

    \section{Busca em Largura (BFS)}

    A busca em largura explora o grafo nível por nível, visitando todos os vizinhos a
    distância $k$ antes de explorar distância $k+1$.

    \textbf{Características:}
    \begin{itemize}
        \item Usa uma \textbf{fila} (FIFO)
        \item Complexidade: $O(V + E)$
        \item \textbf{Garante caminho com menor número de arestas}
        \item Marca visitados ao \textit{enfileirar}, não ao desenfileirar
    \end{itemize}

    \textbf{Vantagem crítica:} em grafos com ciclos, marcar visitados ao enfileirar previne
    que o mesmo vértice seja adicionado à fila múltiplas vezes.

    ---

    \section{Algoritmo de Dijkstra}

    O algoritmo de Dijkstra calcula o caminho de menor custo entre um nó de origem e todos
    os demais nós em um grafo com pesos não-negativos.

    \textbf{Características:}
    \begin{itemize}
        \item Usa \textbf{fila de prioridade} (min-heap)
        \item Complexidade: $O(E \log V)$ com heap binário
        \item Retorna custo mínimo acumulado
    \end{itemize}

    \subsection{Adaptação: Custo Inverso para Maximizar Dependência}

    Neste projeto, adaptamos Dijkstra para encontrar o caminho de \textbf{máxima dependência}
    de comunicação. A intuição é:

    \begin{itemize}
        \item Peso alto = muitos e-mails = relação forte
        \item Para favorecer arestas pesadas, convertemos: $\text{custo} = \frac{1}{\text{peso}}$
        \item Dijkstra minimiza custo $\Rightarrow$ prefere arestas pesadas
    \end{itemize}

    \textbf{Exemplo:}
    \begin{itemize}
        \item Aresta com peso 10: custo = 0.1
        \item Aresta com peso 1: custo = 1.0
        \item Dijkstra minimiza soma de custos $\Rightarrow$ escolhe aresta com peso 10
    \end{itemize}

    ---

    \section{Tratamento de Ciclos}

    Todos os algoritmos implementados tratam ciclos explicitamente:

    \begin{itemize}
        \item \textbf{DFS}: conjunto \texttt{visited} previne revisitas
        \item \textbf{BFS}: marca visitados ao enfileirar
        \item \textbf{Dijkstra}: conjunto \texttt{settled} garante cada vértice é processado uma vez
    \end{itemize}

% ============================================================================
% METODOLOGIA E IMPLEMENTAÇÃO
% ============================================================================

    \chapter{Metodologia e Implementação}

    \section{Arquitetura em Camadas}

    O projeto foi organizado em 7 pacotes (camadas):

    \begin{enumerate}
        \item \textbf{model}: classes de dados puros (Vertex, Edge, PathResult)
        \item \textbf{graph}: estrutura do grafo (ContactGraph)
        \item \textbf{service}: algoritmos (DFS, BFS, DistanceCalculator, CriticalPathService)
        \item \textbf{parser}: leitura e parsing do dataset Enron
        \item \textbf{util}: utilitários (validação de e-mail, fontes)
        \item \textbf{design}: sistema de temas (cores, dark/light)
        \item \textbf{view}: interface gráfica (Swing + GraphStream)
    \end{enumerate}

    Esta separação segue os princípios SOLID:

    \begin{itemize}
        \item \textbf{SRP} (Single Responsibility): cada classe tem uma responsabilidade
        \item \textbf{OCP} (Open/Closed): extensível sem modificar código existente
        \item \textbf{DIP} (Dependency Inversion): depende de abstrações, não de implementações
    \end{itemize}

    ---

    \section{Estrutura de Dados: Grafo com Lista de Adjacência}

    \begin{lstlisting}[language=Java]
public class ContactGraph implements Serializable {
    private final Map<String, Vertex> vertexIndex;  // email -> Vertex
    private final Map<Vertex, Map<Vertex, Edge>> adjacency;  // origem -> (destino -> Edge)
    private final Set<String> ownerEmails;  // senders (150 usu\'arios Enron)
}
    \end{lstlisting}

    \textbf{Operações:}

    \begin{itemize}
        \item \texttt{addVertex(email)}: $O(1)$ médio (HashMap.computeIfAbsent)
        \item \texttt{addEdge(from, to)}: $O(1)$ médio; incrementa peso se existe
        \item \texttt{outDegree(v)}: $O(1)$ (tamanho do mapa interno)
        \item \texttt{inDegree(v)}: $O(V + E)$ (scan necessário)
        \item \texttt{vertexCount()}: $O(1)$
        \item \texttt{edgeCount()}: $O(V)$ (soma de tamanhos de mapas internos)
    \end{itemize}

    ---

    \section{Algoritmos Implementados}

    \subsection{1. Busca em Profundidade (DFS)}

    \textbf{Classe:} \texttt{DepthFirstSearch.java}

    \textbf{Método principal:}

    \begin{lstlisting}[language=Java]
public PathResult search(ContactGraph graph, String originEmail, String destinationEmail)
    \end{lstlisting}

    \textbf{Detalhes de implementação:}

    \begin{itemize}
        \item Usa \texttt{ArrayDeque} como pilha (evita StackOverflowError)
        \item Mantém conjunto \texttt{visited} para ciclos
        \item Mapa \texttt{predecessor} para reconstrução de caminho
        \item Retorna \texttt{PathResult} com lista de vértices ordenados
    \end{itemize}

    \textbf{Garantias:}
    \begin{itemize}
        \item Encontra caminho se existir (completo)
        \item Não necessariamente o mais curto (unlike BFS)
        \item Sem loops infinitos em ciclos
    \end{itemize}

    ---

    \subsection{2. Busca em Largura (BFS)}

    \textbf{Classe:} \texttt{BreadthFirstSearch.java}

    \textbf{Método principal:}

    \begin{lstlisting}[language=Java]
public PathResult search(ContactGraph graph, String originEmail, String destinationEmail)
    \end{lstlisting}

    \textbf{Detalhe crítico:} marca visitados ao \textbf{enfileirar}, não ao desenfileirar.

    \begin{lstlisting}[language=Java]
queue.add(source);
visited.add(source);  // <- Marca AQUI, nao ao dequeue
predecessor.put(source, null);

while (!queue.isEmpty()) {
    Vertex current = queue.poll();
    // ... processar
    for (Edge edge : graph.getOutEdges(current)) {
        Vertex neighbour = edge.getDestination();
        if (!visited.contains(neighbour)) {
            visited.add(neighbour);     // <- Marca ao enfileirar
            predecessor.put(neighbour, current);
            queue.add(neighbour);
        }
    }
}
    \end{lstlisting}

    \textbf{Garantia:} retorna caminho com \textbf{menor número de arestas}.

    ---

    \subsection{3. Distância D (BFS por Níveis)}

    \textbf{Classe:} \texttt{DistanceCalculator.java}

    \textbf{Método principal:}

    \begin{lstlisting}[language=Java]
public List<Vertex> getVerticesAtDistance(ContactGraph graph, String originEmail, int distance)
    \end{lstlisting}

    \textbf{Lógica:}

    \begin{itemize}
        \item Mantém lista de vértices em ``nível atual''
        \item Para cada iteração de 1 a \texttt{distance}:
        \begin{itemize}
            \item Expande todos vizinhos de saída do nível atual
            \item Adiciona ao nível seguinte (se não visitado)
        \end{itemize}
        \item Retorna apenas vértices do nível exato \texttt{distance}
    \end{itemize}

    \textbf{Casos especiais:}
    \begin{itemize}
        \item $D = 0$: retorna apenas o nó de origem
        \item Grafo esgotado antes de $D$: retorna lista vazia
    \end{itemize}

    ---

    \subsection{4. Caminho Crítico (Dijkstra Adaptado)}

    \textbf{Classe:} \texttt{CriticalPathService.java}

    \textbf{Método principal:}

    \begin{lstlisting}[language=Java]
public PathResult computeCriticalPath(ContactGraph graph, String originEmail, String destinationEmail)
    \end{lstlisting}

    \textbf{Adaptação chave:}

    \begin{lstlisting}[language=Java]
// custo inverso favorece arestas pesadas
double newCost = dist.get(current) + edge.getInverseCost();  // 1.0 / weight
    \end{lstlisting}

    \textbf{Saída:}

    \begin{itemize}
        \item Lista de vértices no caminho
        \item Custo inverso total (métrica de Dijkstra)
        \item Dependência acumulada (soma dos pesos originais)
    \end{itemize}

    ---

    \section{Leitura e Parsing do Dataset}

    \textbf{Classe:} \texttt{EnronDatasetReader.java}

    \textbf{Fluxo:}

    \begin{enumerate}
        \item Escaneia diretório \texttt{maildir/} (150 usuários)
        \item Para cada usuário, lê subpastas \texttt{sent/} e \texttt{\_sent\_mail/}
        \item \textbf{Deduplicação}: calcula SHA-256 de cada arquivo
        \begin{itemize}
            \item Se mesmo hash em \texttt{sent/} e \texttt{\_sent\_mail/}: processa apenas uma vez
            \item Evita inflate artificial de pesos
        \end{itemize}
        \item Para cada arquivo único:
        \begin{itemize}
            \item Parse com \texttt{EmailParser} (RFC 2822)
            \item Extrai remetente (From:) e destinatários (To:)
            \item Normaliza e-mails (lowercase, trim)
            \item Adiciona aresta para cada destinatário
        \end{itemize}
        \item Serializa grafo em \texttt{graph.bin} para cache
    \end{enumerate}

    \textbf{Performance:}

    \begin{itemize}
        \item Primeira execução: 1-3 minutos (scan completo + parse)
        \item Execuções subsequentes: < 1 segundo (carrega cache)
        \item Flag \texttt{--rebuild} força re-parse
    \end{itemize}

% ============================================================================
% APRENDIZADOS E CONCLUSÕES
% ============================================================================

    \chapter{Aprendizados e Conclusões}

    \section{Conceitos Reforçados}

    \subsection{Estruturas de Dados}

    \begin{itemize}
        \item \textbf{Maps e Sets}: eficiência de lookup e membership testing
        \item \textbf{Listas vs. Deques}: LIFO (pilha) vs. FIFO (fila) e suas implicações
        \item \textbf{Filas de prioridade}: heap binário para Dijkstra eficiente
        \item \textbf{Representação de grafos}: adjacência vs. matriz vs. edge list
    \end{itemize}

    \subsection{Algoritmos de Grafos}

    \begin{itemize}
        \item \textbf{DFS e BFS}: diferenças semânticas e quando usar cada uma
        \item \textbf{Tratamento de ciclos}: importância de marcar visitados
        \item \textbf{Dijkstra}: otimização com fila de prioridade
        \item \textbf{Adaptações:} como modificar algoritmos para problemas específicos (custo inverso)
    \end{itemize}

    \subsection{Engenharia de Software}

    \begin{itemize}
        \item \textbf{SRP:} código mais testável e manutenível
        \item \textbf{Documentação:} JavaDoc não é \"overhead'', é comunicação
        \item \textbf{Cache:} serialização para performance (graph.bin)
        \item \textbf{Tratamento de erros:} validação robusta de entradas
    \end{itemize}

    ---

    \section{Desafios Enfrentados}

    \begin{enumerate}
        \item \textbf{Deduplicação}: como identificar e-mails duplicados entre pastas?
        \textbf{Solução:} SHA-256 hashing.

        \item \textbf{Performance}: parsing de 600 mil arquivos leva tempo.
        \textbf{Solução:} cache binário com Java Serialization.

        \item \textbf{Dijkstra adaptado}: como inverter a semântica de custo?
        \textbf{Solução:} aplicar transformação $\text{custo} = 1/\text{peso}$.

        \item \textbf{Ciclos}: evitar loops infinitos em grafos com ciclos.
        \textbf{Solução:} marcar visitados antes de explorar vizinhos.
    \end{enumerate}

    ---

    \section{Possíveis Extensões}

    \begin{itemize}
        \item Implementar detecção de comunidades (ex: algoritmo de Louvain)
        \item Adicionar análise de centralidade (betweenness, closeness)
        \item Suportar grafos não-direcionados
        \item Exportar grafo em formato GraphML ou Gephi
        \item Paralelizar parsing com Java Stream parallelism
        \item Adicionar testes unitários (JUnit)
    \end{itemize}

    ---

    \section{Conclusão}

    Este projeto foi uma aplicação prática e abrangente dos conceitos de estruturas de dados
    e algoritmos de grafos. Através da implementação, do zero, de DFS, BFS, cálculo de distância
    e Dijkstra adaptado, consolidei entendimento profundo sobre:

    \begin{itemize}
        \item Como diferentes estratégias de traversal afetam o resultado
        \item Importância de estruturas de dados eficientes (map vs. array)
        \item Trade-offs entre espaço e tempo
        \item Aplicação de conceitos teóricos em problemas reais
    \end{itemize}

    A arquitetura em camadas e a aderência aos princípios SOLID demonstraram que código bem
    estruturado é essencial não apenas para manutenibilidade, mas também para facilitar testes
    e validação.

    ---



% ============================================================================
% APÊNDICES
% ============================================================================

    \appendix

    \chapter{Diagramas UML}
    \setcounter{figure}{0}

    \diagramWideLandscapePage{Diagrama de Classes}{Enron_ContactGraph_Classes.png}{Diagrama de Classes}
    \diagramLandscapePage{Diagrama de Caso de Uso}{Caso_de_Uso-Diagrama_de_Caso_de_Uso__Enron_Graph_Analyzer.png}{Diagrama de Caso de Uso}
    \diagramPortraitPage{Sequencia de Construcao do Grafo}{Sequencia_Construcao_Grafo-Graph_Construction_Sequence.png}{Sequencia de Construcao do Grafo}
    \diagramPortraitPage{Sequencia de Busca DFS e BFS}{Sequencia_Busca-DFS_and_BFS_Search_Sequence.png}{Sequencia de Busca DFS e BFS}
    \diagramPortraitPage{Sequencia de Caminho Critico}{Sequencia_Caminho_Critico-Critical_Path_Sequence__Adapted_Dijkstra_.png}{Sequencia de Caminho Critico}
    \diagramLandscapePage{Diagrama de Componentes}{Diagrama_Componentes-Diagrama_de_Componentes__Enron_Graph_Analyzer.png}{Diagrama de Componentes}
    \diagramPortraitPage{Diagrama de Dominio}{Diagrama_Dominio-Diagrama_de_Dominio__Enron_Graph_Analyzer.png}{Diagrama de Dominio}
    \diagramPortraitPage{Diagrama de Atividade}{Diagrama_Atividade-Diagrama_de_Atividade__Fluxo_Principal_do_Enron_Graph_Analyzer.png}{Diagrama de Atividade}
    \diagramPortraitPage{Diagrama de Estado}{Diagrama_Estado-Diagrama_de_Estado__Enron_Graph_Analyzer.png}{Diagrama de Estado}
    \diagramLandscapePage{Diagrama de Comunicacao}{Diagrama_Comunicacao-Diagrama_de_Comunicacao__Interacao_Principal.png}{Diagrama de Comunicacao}

    \chapter{Exemplos de Saída}

    \section{Modo Demonstração}

    \begin{lstlisting}[language=bash]
$ mvn exec:java

=== DEMO MODE ===

--- 1. GRAPH STATISTICS ---
  Vertices : 5
  Edges    : 6

--- 2. TOP 20 OUT-DEGREE (most active senders) ---
   1. alice@company.com (degree=2)
   2. bob@company.com (degree=2)
   3. carol@company.com (degree=1)
   4. dave@company.com (degree=1)

--- 3. TOP 20 IN-DEGREE (most contacted recipients) ---
   1. dave@company.com (degree=2)
   2. eve@company.com (degree=2)
   3. bob@company.com (degree=1)
   4. carol@company.com (degree=1)

--- 4-7. INTERACTIVE SEARCH PANEL ---
  Use the GUI windows to run DFS, BFS, Distance D and Critical Path.

[Janela Swing é exibida]
    \end{lstlisting}

    ---

    \section{Consulta BFS no Demo}

    \begin{lstlisting}
Origem: alice@company.com
Destino: eve@company.com
Algoritmo: BFS

RESULTADO:
alice@company.com -> bob@company.com -> eve@company.com

Comprimento: 2 arestas
Custo total das arestas: 2 mensagens (bob->eve tem peso 1)
    \end{lstlisting}

    ---

\end{document}
